%% sections/05-discussion.tex

\section{Discussion}
\label{sec:discussion}

\subsection{Why Character Logic Produces Better Payoffs Than DM Scripting}

The 90\% session coverage rate and 100\% memorable-NPC-moment rate for
sessions with arc-completions together support a claim that goes beyond
coverage: the \textit{origin} of an NPC moment --- whether it was earned
by the party or scripted by the DM --- is a meaningful quality axis, not
merely a design preference.

There is a structural account of why this is so. When a party earns an NPC
moment through specific, observable action --- when the party's behavior is
the direct cause of the NPC's arc-completion --- the moment belongs to the
session in a way that DM-scripted moments cannot. Players may not articulate
this distinction, but its effects are visible in panel review language: arc-
completion moments are described as ``the session''s best moment,'' as things
that ``the party did,'' as moments that ``had been building.'' Improvised NPC
beats --- even effective ones --- are described as ``what the DM gave us.''
The arc-completion pattern changes the causal attribution of emotional payoffs
from DM-to-players to players-to-session.

This reframes the pattern's primary design implication. The intervention is
not ``tell DMs to time NPC moments better'' or ``tell DMs to be more
emotionally attentive.'' It is ``tell adventure designers to specify, before
play, what each NPC is capable of saying and exactly what the party must do
to see it.'' The constraint forces articulation that improvisation avoids: the
designer must answer the question of what the NPC can offer and under what
conditions, which is a harder design question than ``what would this NPC say
if the party were friendly with them.''

\subsection{The Information-Holder NPC Variant}

The five exemplar NPCs in Section~\ref{sec:methodology} include one
information-holder type (Aldras Bent), and the Campaign 5 discussion in
Section~\ref{sec:evaluation} extends this with Bessa Tor and Pell Dorvast.
The information-holder variant warrants specific attention because it is
the type most likely to be confused with interrogation design.

In conventional adventure design, NPCs who hold information are designed as
interrogation targets: the party must find leverage, make skill checks, or
build rapport to extract the information. The firing condition in these
designs is typically a threshold (persuasion DC 16, or three successful
interaction beats) rather than a specific observable action. This design
structure makes the NPC's information contingent on party capability, which
is mechanically tractable but narratively flat.

The information-holder arc-completion inverts this. The firing condition is
not a threshold but a behavioral signature: the party must demonstrate a
specific form of professional respect that the NPC's arc was designed to
detect. Aldras Bent does not name his handler because the party passed a
persuasion check; he names him because the party showed him the ledger-date
and waited, which is the specific gesture of professional acknowledgment his
arc-completion condition required. The information is the arc-completion act,
not the reward for passing a check.

This distinction has practical implications for adventure design: information-
holder NPCs should be designed with arc-completion conditions, not interrogation
thresholds, when the goal is narrative memorability rather than information
delivery efficiency. An NPC who voluntarily offers information under conditions
of professional respect produces a more memorable moment than an NPC who
surrenders information under persuasion pressure, even when the information
transferred is identical.

\subsection{The Bilateral Arc-Completion}

Four sessions in the corpus featured bilateral arc-completions: scenes in which
both an NPC and a PC fired arc-completion moments simultaneously, each creating
the condition for the other. These instances (Sera+Morreth in C2-S03,
Thessaly+Vorn in C2-S05, Asholt+Saren in C4-S07, and one instance in C1-S07)
produced the four highest single-scene Atmospheric Landing scores in the corpus.

Bilateral arc-completions require a specific design condition: both the NPC's
arc-completion and the PC's arc-completion must share a trigger domain --- a
subject matter or action space where each party's conditions can be met by the
same gesture. In Sera+Morreth, the shared domain was naming: Sera named the
settlement she had been protecting; Morreth named Ven, the person he had lost
protecting it. The party action (Sera's naming) met both the PC's own arc-
completion condition and Morreth's NPC condition simultaneously.

Bilateral arc-completions cannot be forced through design alone; they require
both design (shared trigger domain) and emergence (the party's behavior must
actually converge on the shared space). The four instances in the corpus were
all noted as unplanned in the sense that no session was designed with a
bilateral arc-completion as a target; they emerged from the intersection of
NPC and PC arc designs. This suggests a design principle: ensure NPC and PC
arc designs share trigger domains where possible, then allow bilateral
convergence to emerge from party behavior. Do not pre-script the bilateral.

\subsection{Arc-Completion and Drama Management as Complementary Approaches}

The arc-completion pattern is most usefully understood as complementary to drama
management rather than competing with it. Drama management, as exemplified by
\citet{mateas2003faade}, operates at the session level: it monitors emergent play
and steers narrative toward designed endpoints through real-time intervention.
Arc-completion operates at the NPC level: it specifies, before play, the conditions
under which a character will say or do something emotionally significant. These
are different levels of the design stack, and they are compatible.

A system combining pre-specified arc-completion conditions with a lightweight drama
manager for pacing could achieve both reliability and adaptability. The arc-completion
layer would guarantee that NPC emotional moments fire on authentic terms (condition
met, character logic fires); the drama management layer would handle session-level
concerns such as encounter sequencing, pacing adjustment, and response to unexpected
player choices that do not engage any NPC's arc-completion conditions. In the Marathon
corpus, the session-runner functions as a minimal drama manager for these pacing
concerns, while the NPC file system provides the arc-completion layer. The two
functions do not interfere; they operate on different parts of the session's design
space. This separation suggests a practical design principle: let drama management
handle session shape; let arc-completion design handle NPC authenticity. Do not ask
drama management to time emotional moments, and do not ask arc-completion to recover
from session-level pacing failures.

\subsection{Limits and Threats to Validity}

\paragraph{Single AI system.}
The corpus is generated by a single AI system (Marathon's session-runner)
running all DM functions. This system enforces firing-condition checking
mechanically: the AI cannot script an arc-completion without checking the
condition. The zero DM-scripted moment count in the corpus is therefore partly
a system property, not purely a design property. In human-run TRPG, the
pattern requires DM discipline to hold. We cannot report data on whether human
DMs maintain this discipline; future work should apply the pattern in human-run
sessions and measure compliance with DM discipline notes.

\paragraph{No player studies.}
All quality assessments in the corpus are produced by AI-simulated panel
reviewers using a calibrated rubric, not by human players. The claim that
character-logic arc-completions produce better player experiences than DM-
scripted moments is an inference from panel scoring patterns, not a direct
measurement of player affect. Future work should instrument the pattern in
human-played sessions with player-experience measures (post-session recall
interviews, affect ratings, or session-satisfaction surveys).

\paragraph{Thematic range.}
Despite seven campaign archetypes, the corpus does not include campaigns with
antagonist-only parties, sandbox exploration structures, or player-driven
emergent narrative (e.g., Powered by the Apocalypse format). The arc-completion
pattern's core mechanism (specify the condition and the text before play; let
character logic fire it) should transfer, but the specific design template may
require adaptation for systems where NPC behavior is defined by player-facing
moves rather than DM-administered arc files.

\paragraph{Correlation vs. causation.}
The strong correlation ($r = 0.89$) between arc-completion presence and
session-memorable NPC moments is not proof of causation. Sessions with well-
designed arc-completions are also sessions with stronger overall NPC design,
more sustained party investment in NPC relationships, and more carefully
structured encounter sequencing. The arc-completion may be a proxy for overall
design quality rather than an independent driver of memorable moments. The
failure analysis provides partial evidence against this interpretation (sessions
with high overall scores but no arc-completions still fail to produce memorable
NPC moments at high rates), but a controlled comparison with matched-quality
session designs would be required to establish causation cleanly.
